\documentclass{article}

\usepackage{amsmath,amssymb}
\usepackage{xurl}
\usepackage[colorlinks=true,linkcolor=blue,citecolor=blue,urlcolor=blue]{hyperref}

% Shared commands for notes in docs/paper-gaps/.

\DeclareUnicodeCharacter{2080}{\ifmmode{}_{0}\else\textsubscript{0}\fi}
\DeclareUnicodeCharacter{2081}{\ifmmode{}_{1}\else\textsubscript{1}\fi}
\DeclareUnicodeCharacter{2082}{\ifmmode{}_{2}\else\textsubscript{2}\fi}
\DeclareUnicodeCharacter{2099}{\ifmmode{}_{n}\else\textsubscript{n}\fi}

\newcommand{\Lean}{\textsc{Lean}}
\ExplSyntaxOn
\NewDocumentCommand{\leanid}{m}{
  \ifmmode
    \textsf{\detokenize{#1}}
  \else
    \group_begin:
    \urlstyle{sf}
    \tl_set:Nn \l_tmpa_tl {#1}
    \tl_replace_all:Nnn \l_tmpa_tl {\_} {_}
    \exp_args:NV \path \l_tmpa_tl
    \group_end:
  \fi
}
\ExplSyntaxOff
\providecommand{\tr}{\operatorname{tr}}
\newcommand{\tnleanIssueBase}{https://github.com/LionSR/TNLean/issues/}
\newcommand{\tnleanPrBase}{https://github.com/LionSR/TNLean/pull/}
\newcommand{\tnleanDocBase}{https://sirui-lu.com/TNLean/docs/}
\newcommand{\ghissue}[1]{\href{\tnleanIssueBase#1}{GitHub issue \##1}}
\newcommand{\ghpr}[1]{\href{\tnleanPrBase#1}{PR \##1}}
\newcommand{\doclink}[1]{\href{\tnleanDocBase#1.html}{\path{#1}}}

% Dirac notation and the identity operator, as in blueprint/src/macros/common.tex.
% \providecommand keeps note-local definitions authoritative.
\usepackage{dsfont}
\providecommand{\ket}[1]{|#1\rangle}
\providecommand{\bra}[1]{\langle #1|}
\providecommand{\braket}[2]{\langle #1 | #2 \rangle}
\providecommand{\ketbra}[2]{|#1\rangle\!\langle #2|}
\providecommand{\Id}{\mathds{1}}
\providecommand{\one}{\mathds{1}}
\providecommand{\C}{\mathbb{C}}
\providecommand{\MN}[1]{M_{#1}(\C)}
\providecommand{\MD}{M_D(\C)}

% External referencing for paper sources.  The build helper writes sanitized
% auxiliary files under build/paper-aux/ and excludes duplicated source labels.
\usepackage{xr}

% Import a generated paper auxiliary file with a caller-supplied label prefix.
% The candidate paths support builds from docs/paper-gaps/, the repository root,
% and build/paper-aux/ itself.
\newcommand{\importpaperaux}[2]{%
  \IfFileExists{../../build/paper-aux/#2.aux}{%
    \externaldocument[#1]{../../build/paper-aux/#2}%
  }{%
    \IfFileExists{build/paper-aux/#2.aux}{%
      \externaldocument[#1]{build/paper-aux/#2}%
    }{%
      \IfFileExists{#2.aux}{%
        \externaldocument[#1]{#2}%
      }{}%
    }%
  }%
}

% General external reference macros
\newcommand{\paperref}[2]{\ref{#1-#2}}
\newcommand{\papereqref}[2]{(\ref{#1-#2})}

% CPSV16 source (1606.00608 / MPDO-22-12-17-2.tex) external document and shortcuts
\importpaperaux{cpsv16-}{1606.00608/MPDO-22-12-17-2-tnlean}
\newcommand{\cpsvref}[1]{\paperref{cpsv16}{#1}}
\newcommand{\cpsveqref}[1]{\papereqref{cpsv16}{#1}}

% Verdict marker: \gapnote{<kind>}{<status>}. Kind is the policy
% classification (clarification, local-correction, scope-restriction,
% unfaithful, false-source, open-gap); status is open, wip (elimination
% underway), resolved, or historical. Machine-read by the site index;
% prints a verdict line.
\newcommand{\gapnote}[2]{\par\noindent\textbf{Verdict:} #1 (#2).\par}


\title{Removal Note: the strict weight-modulus route to the\\
       MPS Fundamental Theorem}
\author{}
\date{2026-06-01}

\begin{document}
\maketitle

This note records, before its deletion, the strict weight-modulus
(``dominant-block peeling'') route that the formal development used for the
non-periodic Matrix Product State Fundamental Theorem until it was superseded by
the sector-coefficient comparison route of
\cite[\S II]{Cirac2016MPDO_arXiv}.\footnote{%
  The replacement route is the one anchored in
  \texttt{blueprint/src/chapter/ch10\_bnt.tex} and
  \texttt{ch11\_fundamental\_theorem\_proof.tex}; its Lean home is
  \doclink{TNLean/MPS/FundamentalTheorem/SectorBNT/Fundamental}.}
The purpose of the note is to keep a coherent mathematical record of what the
route assumed, what it was designed to prove, why the assumption is not
faithful to the source canonical form, and why nothing load-bearing is lost when
the route is removed. It complements
\path{docs/paper-gaps/cpsv16_nondominant_per_block_projection.tex}, which records
the analytic failure of the projection step in detail.

\section{The strict canonical form and the predicates being removed}

Fix a finite family of normal MPS site tensors $(A_k)_{k=0}^{r-1}$ with nonzero
complex weights $\mu_k$, assembled into the block-diagonal tensor
$A=\bigoplus_k \mu_k A_k$. Its matrix-product vector at chain length $N$ expands
as
\begin{align}
    V^{(N)}(A) =  \sum_{k=0}^{r-1} \mu_k^{\,N}\, V^{(N)}(A_k).
    \label{eq:assemble}
\end{align}
The removed route strengthened the canonical form by a \emph{strict} ordering of
the weight moduli together with a dominant normalization,
\begin{align}
    |\mu_0| > |\mu_1| > \cdots > |\mu_{r-1}|,
    \qquad |\mu_0| = 1.
    \label{eq:strict}
\end{align}
In the formalization \eqref{eq:strict} was carried by the predicate
\leanid{HasStrictOrderedNonzeroWeights} (field
\leanid{mu_strict_anti}) and by the canonical-form structure
\leanid{IsNormalCanonicalFormBNT} (fields \leanid{mu_strict_anti} and
\leanid{mu_dom_norm_one}).\footnote{%
  \leanid{HasStrictOrderedNonzeroWeights} and
  \leanid{IsNormalCanonicalForm.ofStrictSeparatedData} in
  \doclink{TNLean/PiAlgebra/CanonicalFormSepAux};
  \leanid{MPSTensor.IsNormalCanonicalFormBNT} in
  \doclink{TNLean/MPS/BNT/Construction}.}
The block-separation and ``fundamental theorem for the canonical form''
declarations built on \eqref{eq:strict} lived in
\doclink{TNLean/PiAlgebra/CanonicalFormSep}; the consuming injectivity-after-blocking
results lived in
\doclink{TNLean/MPS/CanonicalForm/SectorComparison/NormalityChain}.

\section{The argument the route was designed to support}

The strict ordering \eqref{eq:strict} was not a cosmetic convention. It was the
engine of a dominant-block peeling proof of the equal-MPV uniqueness statement.
Given two assembled families $A,B$ with $V^{(N)}(A)=V^{(N)}(B)$ for all $N$,
normalize \eqref{eq:assemble} by the dominant weight $\mu_0^N$. Since
$|\mu_k/\mu_0|<1$ for every $k\neq 0$ under \eqref{eq:strict}, the renormalized
contribution of every sub-dominant block tends to zero,
\begin{align}
    \left(\tfrac{\mu_k}{\mu_0}\right)^{\!N} \xrightarrow[N\to\infty]{} 0
    \qquad (k \neq 0),
    \label{eq:decay}
\end{align}
so in the limit only the leading block survives. One then projects onto the
leading matrix-product vector, matches the leading blocks of $A$ and $B$,
subtracts (``peels'') the matched leading block, and recurses on the strictly
shorter remaining family. The strict gap in \eqref{eq:strict} supplies both the
unique dominant block at each step and the spectral separation that makes
\eqref{eq:decay} and the induction on the number of blocks terminate. For the
equal-MPV case this is sound, because the leading-pair comparison produces the
\emph{exact} identity $\mu^A_0 = \mu^B_0\,\zeta$ with $|\zeta|=1$, after which the
leading summand cancels exactly and the tail again satisfies an equality of the
same shape.

\section{Why strict ordering is not faithful to the source canonical form}

The source canonical form of \cite[Def.~4.2, lines~1864--1884]{Cirac2021Matrix}
groups the normal blocks into sectors. A sector indexed by $j$ may contain
several copies of one normal tensor with weights $\mu_{j,q}$, so the coefficient
attached to a basis-of-normal-tensors vector is the \emph{power sum}
\begin{align}
    c_j(N) =  \sum_q \mu_{j,q}^{\,N},
    \label{eq:powersum}
\end{align}
not a single scalar power. The strict ordering \eqref{eq:strict} silently
collapses \eqref{eq:powersum} to one term $\mu_j^N$, i.e.\ it assumes one copy
per sector, and it forbids equal-modulus blocks. Both restrictions exclude
canonical forms that the source theorem must cover:
\begin{itemize}
    \item $C \oplus (-C)$, with weights $1$ and $-1$ of equal modulus, has
          coefficient $c(N) = 1 + (-1)^N$, which is not a single scalar power;
    \item $C \oplus D$ with two distinct non-gauge-equivalent normal tensors and
          equal weights $(1,1)$ has two basis-of-normal-tensors vectors of equal
          modulus;
    \item $C \oplus \tfrac12 C$ is one sector with two copies of unequal modulus,
          coefficient $1 + (\tfrac12)^N$.
\end{itemize}
Grouping into a basis of normal tensors is governed by gauge-phase equivalence of
the block tensors, not by distinctness of weight moduli; equal-modulus sectors
are part of the theorem, not an afterthought.

\section{Why the peeling argument fails beyond the equal-MPV case}

For the proportional statement, where $V^{(N)}(A)=c_N\,V^{(N)}(B)$ with a nonzero
scalar sequence $c_N$, the peeling argument breaks even under \eqref{eq:strict}.
Projecting onto a \emph{non-dominant} block $B_{k_0}$ with $k_0\neq 0$ sends both
sides to zero: the left side is a bounded coefficient times a decaying overlap,
and the diagonal right-hand term carries the factor $|\mu_{k_0}/\mu_0|^N \to 0$.
No contradiction is available, and the exact leading-coefficient identity that the
recursion needs is false at finite chain length. This obstruction, and the two
independent probes confirming it, are recorded in
\path{docs/paper-gaps/cpsv16_nondominant_per_block_projection.tex}.\footnote{%
  See also \path{audits/2026-05-13_cpsv16_ft_discharge_attempt.md} and
  \path{audits/2026-05-13_cpsv16_ft_exact_leading_coeff.md}.}

\section{The replacement and the verdict}

The current route does not take limits. At a fixed sufficiently large $N$ it
isolates $B_{k_0}$, observes that the remaining terms lie in the span of
$\{V^{(N)}(A_j)\}_j \cup \{V^{(N)}(B_k) : k\neq k_0\}$, and uses eventual linear
independence of the full combined family to force the isolated coefficient
$c_N\mu_{k_0}^N$ to vanish, contradicting $c_N\neq 0$ and
$\mu_{k_0}\neq 0$. This identifies coefficients even when they decay, and it
needs no ordering of the moduli. It is realized over the raw sector weights and
power-sum coefficients \eqref{eq:powersum} of \leanid{SectorDecomposition}, with
the combined-family statement
\leanid{MPSTensor.IsBNTCanonicalForm.combined_family_eventually_li} and the
non-decay of unit-modulus power sums
\leanid{MPSTensor.sum_pow_not_tendsto_zero_of_unit_modulus}.

The verdict is that the strict route was genuinely load-bearing for the old
peeling architecture, but that architecture was itself the one-copy specialization
just described. The headline Fundamental-Theorem declarations
\leanid{MPSTensor.ft_sector_bnt_equal_global_gauge} and
\leanid{MPSTensor.ft_sector_bnt_equal_mps_gaugeEquiv_literal} do not depend, by
proof term, on \eqref{eq:strict} or on any declaration of
\doclink{TNLean/PiAlgebra/CanonicalFormSep}. Removing the strict route therefore
deletes a parallel, strictly weaker development; it does not weaken the formalized
theorem.

\bibliographystyle{alpha}
\bibliography{references}

\end{document}